\documentclass{ieeeaccess}

\usepackage{cite}
\usepackage{amsmath,amssymb,amsfonts}
\usepackage{graphicx}
\usepackage{textcomp}
\usepackage{booktabs}
\usepackage{tabularx}
\usepackage{hyperref}
\usepackage{multirow}
\usepackage{bm}

\makeatletter
\AtBeginDocument{\DeclareMathVersion{bold}
\SetSymbolFont{operators}{bold}{T1}{times}{b}{n}
\SetSymbolFont{NewLetters}{bold}{T1}{times}{b}{it}
\SetMathAlphabet{\mathrm}{bold}{T1}{times}{b}{n}
\SetMathAlphabet{\mathit}{bold}{T1}{times}{b}{it}
\SetMathAlphabet{\mathbf}{bold}{T1}{times}{b}{n}
\SetMathAlphabet{\mathtt}{bold}{OT1}{pcr}{b}{n}
\SetSymbolFont{symbols}{bold}{OMS}{cmsy}{b}{n}
\renewcommand\boldmath{\@nomath\boldmath\mathversion{bold}}}
\makeatother

\hyphenation{pan-demic pre-pared-ness fore-cast-ing mor-tal-i-ty
            epi-de-mi-ol-og-i-cal com-pen-sa-tion}

\begin{document}
\history{Date of publication xxxx 00, 0000, date of current version xxxx 00, 0000.}
\doi{10.1109/ACCESS.2024.0000000}

\title{The Illusion of Learning: Provenance-Separated Evaluation of ML Pandemic Forecasting Systems}

\author{\uppercase{D. Author}\authorrefmark{1},
\uppercase{Faculty Mentor}\authorrefmark{2},
\IEEEmembership{Senior Member, IEEE},
\uppercase{Co-Author}\authorrefmark{1},
AND \uppercase{Another Author}\authorrefmark{1}}

\address[1]{Department of Computer Science and Engineering,
Vellore Institute of Technology, Chennai, India
(e-mail: d.author@vit.ac.in)}
\address[2]{Department of Computer Science and Engineering,
Vellore Institute of Technology, Chennai, India
(e-mail: faculty.mentor@vit.ac.in)}

\tfootnote{This work was supported by Vellore Institute of Technology,
Chennai. Corresponding author: Faculty Mentor.}

\markboth
{Author \headeretal: THE ILLUSION OF LEARNING}
{Author \headeretal: THE ILLUSION OF LEARNING}

\corresp{Corresponding author: Faculty Mentor (e-mail: faculty.mentor@vit.ac.in).}

\begin{abstract}
Pandemic ML systems are frequently trained on synthetic or imputed data when real surveillance records are unavailable, yet standard evaluations rarely separate performance by data provenance. We introduce provenance-separated evaluation, a methodology that explicitly tracks data lineage (real vs.\ synthetic) and independently measures model performance across each source. Applied to a multi-model forecasting system of nine specialized models covering seven diseases across 30 Indian states, our audit reveals that six of nine models do not learn epidemiological dynamics in the way their architecture suggests. They either reconstruct deterministic formulas from their own synthetic training targets, function as trivial autoregressive processes, or have negative predictive skill. Even the "real" data pipeline injects synthetic noise (R0 values are uniform random draws). We document a training/inference feature mismatch in the deployed scenario predictor, and show that no production monitoring exists to detect when models fail on real-world data. Our central finding is a 3$\times$ MAPE degradation on real versus synthetic death records (42.1\% vs.\ 14.4\%) concealed by aggregate metrics. We argue that provenance-separated evaluation, model-by-model audits, and pipeline transparency should be standard practice for health ML systems, and release our full evaluation framework and audit trail.
\end{abstract}

\begin{keywords}
Pandemic forecasting, synthetic data, provenance evaluation, machine learning, model audit, negative results, formula reconstruction.
\end{keywords}

\titlepgskip=-21pt

\maketitle

\section{Introduction}
\label{sec:intro}

\PARstart{M}{achine} learning for pandemic preparedness has attracted significant attention following the COVID-19 pandemic~\cite{wang2020,yan2023,hassan2022}. A critical but underappreciated challenge is the extreme scarcity of granular, real-world surveillance records for emerging pathogens. Researchers and practitioners necessarily rely on synthetic or imputed data---generated from epidemiological parameters, mechanistic models, or data-augmentation techniques---to train ML models for forecasting, resource allocation, and risk assessment.

Despite this reliance on non-real data, evaluations of pandemic ML systems rarely separate performance metrics by data provenance. When real and synthetic test records are mixed into a single aggregate score, models that overfit to deterministic generation rules of synthetic data can appear highly performant, masking failure on actual outbreaks. This creates a dangerous feedback loop: a system that impresses on paper may mislead during a real health emergency.

We built a nine-model ML forecasting system for 30 Indian states and seven diseases, then conducted an exhaustive audit. The audit revealed that six of nine models do not learn in the way their reported metrics suggest. Four reconstruct deterministic formulas from synthetic targets. One has negative R\textsuperscript{2}. One has a critical training/inference feature mismatch. Only two models (forecast and mortality) represent genuine ML, and the forecast model reduces to an AR(1) process with 80.5\% of its power from a single lag feature.

This paper documents these findings as a negative result for the health ML community. Our contribution is not the system itself but the evaluation methodology and the cautionary case study it enables.

\textbf{Contributions:}
\begin{enumerate}
    \item \textbf{Provenance-Separated Evaluation}: A framework for tracking data lineage and stratifying model performance by real vs.\ synthetic source. Applied to our system, it reveals a 3$\times$ MAPE degradation on real death records hidden by aggregate scores.
    \item \textbf{Systematic Model Audit}: We classify each of nine models by what it actually learns---formula reconstruction, AR(1) process, or genuine ML---revealing that reported metrics for six models are misleading indicators of predictive skill.
    \item \textbf{AR(1) Dominance}: The forecast model derives 80.5\% of its predictive power from a single lagged feature, functioning as an autoregressive process with 16 extraneous inputs.
    \item \textbf{Training/Inference Mismatch Discovery}: The deployed scenario predictor uses different feature names (current-year vs.\ previous-year disease parameters) than training, compromising predictions silently.
    \item \textbf{Open-Source Audit Framework}: We release our evaluation toolkit, ablation framework, and full audit trail for independent verification.
\end{enumerate}

\section{Related Work}
\label{sec:related}

\subsection{Pandemic Forecasting}
SEIR compartmental models~\cite{kermack1927} remain the epidemiological standard. ML approaches using XGBoost~\cite{chen2016} have shown competitive forecasts~\cite{wang2020,yan2023}. ARIMA~\cite{box2015} and Prophet~\cite{taylor2018} provide interpretable baselines. Our work extends these comparisons by introducing provenance as an evaluation dimension.

\subsection{Synthetic Data in Healthcare ML}
Synthetic data generation is established methodology~\cite{reich2022,gonzales2023}. The CDC recommends simulation-based planning~\cite{cdc2019}. However, fewer than 15\% of health ML papers using synthetic data report provenance-aware evaluation~\cite{chen2023}. Recent work on synthetic data for epidemic forecasting~\cite{marchi2026} argues synthetic data improves model performance---reporting only aggregate metrics. Our paper provides the counterpoint: aggregate metrics on mixed-provenance data can be actively misleading.

\subsection{ML Auditing}
Model cards~\cite{mitchell2019} and pipeline audits~\cite{rajkomar2022} emphasize documentation. Selbst et al.~\cite{selbst2019} identify framing mismatches in applied ML. Schelter et al.~\cite{schelter2023} propose provenance-based debugging. Our work contributes a concrete audit of a health ML pipeline, documenting formula reconstruction, training/inference mismatches, and post-hoc corrections invisible in standard evaluations.

\section{Experimental Platform}
\label{sec:arch}

Our platform is a multi-model backend (FastAPI) with nine specialized models loaded at startup. Table~\ref{tab:model_audit} presents each model alongside its actual behavior as revealed by audit.

\begin{table}[h]
\centering
\caption{Complete Model Audit: Nine Models Classified by Actual Behavior}
\label{tab:model_audit}
\small
\begin{tabular}{lclp{5.5cm}}
\toprule
Model & Algorithm & Actual Behavior & Evidence \\
\midrule
Forecast & XGBoost & AR(1) + 15 extraneous features & 80.5\% importance from lag\_1\_cases \\
Mortality & XGBoost & Genuine ML & R\textsuperscript{2}=0.89, no training script found \\
R0 & XGBoost & Formula reconstruction & R\textsuperscript{2}=0.99 on synthetic formula: $0.96^y (1-0.4v)(1+0.3m)$ \\
Bed & GBR & Negative skill & R\textsuperscript{2}=-0.035, worse than mean prediction \\
Scenario & XGBoost & Overfitting + feature mismatch & 8.3\% train vs.\ 20.1\% test MAPE; inference uses \texttt{avg\_cfr}, trained on \texttt{prev\_avg\_cfr} \\
Patient Risk & Ensemble & Formula reconstruction & 3-model ensemble learns \texttt{compute\_risk\_labels()} formula \\
Lockdown & XGBClassifier & Circular labels & Labels derived from percentile ranks of input features \\
Risk & Deterministic & Pure formula & $\log_{10}$ math, no model loaded \\
Hospital & RF & Hybrid & ML predicts success rate, rankings from CSV \\
\bottomrule
\end{tabular}
\end{table}

\subsection{Model Architecture Summary}
The platform consists of: ForecastPredictor (XGBoost, 16 features, monthly granularity), BedPredictor (GradientBoostingRegressor), MortalityPredictor (XGBoost), ScenarioPredictor (XGBoost, annual total), PatientRiskPredictor (RF+GB+XGB ensemble), HospitalPredictor (RF), R0Predictor (XGBoost), RiskPredictor (deterministic), and LockdownPredictor (XGBoost classifier).

A three-tier fallback chain handles missing data: ML model $\rightarrow$ DB trend $\rightarrow$ hardcoded heuristic.

\subsection{Provenance Tracking}
Every data point carries an \texttt{is\_real} flag and \texttt{source} string persisting through training and evaluation.

\section{Data}
\label{sec:data}

Table~\ref{tab:data} summarizes provenance. Real data constitutes 22\% (3,928 of 17,848 records).

\begin{table}[h]
\centering
\caption{Data Sources and Provenance}
\label{tab:data}
\small
\begin{tabular}{lllr}
\toprule
Disease & Source & Type & Records \\
\midrule
H5N1 & OWID/WHO API (India-filtered) & Real & 352 \\
COVID-19 & COVID19-India API (state-level) & Real & 3,512 \\
SARS 2003 & sars2003.com (post-dedup) & Real & 64 \\
Ebola & Literature-grounded\footnotemark[1] & Synthetic & 3,480 \\
Nipah & Literature-grounded\footnotemark[1] & Synthetic & 3,480 \\
Marburg & Literature-grounded\footnotemark[1] & Synthetic & 3,480 \\
H1N1 & Literature-grounded\footnotemark[1] & Synthetic & 3,480 \\
\bottomrule
\end{tabular}
\end{table}
\footnotetext[1]{WHO CFR/R$_0$ + Indian census population weights}

\subsection{Data Quality Issues}
A deeper audit of the data pipeline reveals critical quality concerns:

\textbf{R0 values are random noise.} In \texttt{real\_data\_ingest.py}, every reproduction rate in the database is generated as \texttt{r0 = np.random.uniform(1.0, 4.0)}. Any model using this column learns patterns from uniform random noise. Similarly, bed demand and ICU demand are generated as \texttt{int(confirmed * np.random.uniform(0.05, 0.15))}---confirming that even the ``real'' data pipeline injects synthetic noise at every step.

\textbf{Provenance labels are inconsistent.} The pandemic service labels only COVID-19 as ``real\_data''; H5N1 (with 352 real OWID records) and SARS (64 real records) are labeled ``who\_parameterized'' alongside purely synthetic diseases.

\textbf{Synthetic data dominates.} Four of seven diseases have zero real surveillance records. Even ``real'' H5N1 data exists only as national aggregates imputed to state level. The boundary between real and synthetic is substantially blurred in practice.

\section{Methodology}
\label{sec:methodology}

\subsection{Evaluation Framework}
Time-series models use chronological 85/15 train/test split with 5-fold TimeSeriesSplit. XGBoost grid: max\_depth $\in \{3,4,5\}$, learning\_rate $\in \{0.01, 0.05, 0.1\}$, reg\_lambda $\in \{1.0, 5.0, 10.0\}$, min\_child\_weight $\in \{3,7\}$, subsample $\in \{0.8, 1.0\}$. We report 95\% confidence intervals via bootstrap ($B=1000$) and Diebold-Mariano tests~\cite{diebold1995}.

\subsection{Provenance-Separated Metrics}
Rather than reporting a single metric across the full test set, we independently measure:
\begin{enumerate}
    \item \textbf{All data}: Standard practice.
    \item \textbf{Real-only}: Records where \texttt{is\_real = True}.
    \item \textbf{Synthetic-only}: Records where \texttt{is\_real = False}.
\end{enumerate}

\subsection{Baseline Models}
We compare against SEIR (national), ARIMA (national and per-state), logistic growth (national), and Ridge regression (per-state).

\subsection{Model Classification Protocol}
We classify each model into one of four categories based on code audit:
\begin{enumerate}
    \item \textbf{Genuine ML}: Model learns from data with no formula-reconstruction or leakage.
    \item \textbf{Formula-reconstruction}: Model learns a deterministic function applied to its own input features.
    \item \textbf{AR(1) dominance}: Model's effective behavior reduces to lag-1 autoregression.
    \item \textbf{Negative skill}: Model performs worse than predicting the mean.
\end{enumerate}

\section{Results}
\label{sec:results}

\subsection{Real vs.\ Synthetic Performance}
\label{sec:real_synth}

Our central finding is the performance gap revealed by provenance-separated evaluation (Table~\ref{tab:real_synth}). On mixed data, XGBoost achieves 11.2\% MAPE for cases and 16.0\% for deaths. Isolating the 465 real-only records exposes a 42\% degradation in case forecasting (15.9\%) and a 3$\times$ degradation in death forecasting (42.1\%).

\begin{table}[h]
\centering
\caption{Forecast Model: Real vs.\ Synthetic Subset}
\label{tab:real_synth}
\small
\begin{tabular}{lccc}
\toprule
Subset & n & MAPE & 95\% CI \\
\midrule
Cases (all) & 3,416 & 11.2\% & [10.8, 11.6] \\
\quad Real only & 465 & 15.9\% & [13.9, 18.3] \\
\quad Synthetic only & 2,951 & 10.5\% & [10.2, 10.8] \\
\addlinespace
Deaths (all) & 3,416 & 16.0\% & [15.2, 16.8] \\
\quad Real only & 465 & 42.1\% & [33.0, 52.2] \\
\quad Synthetic only & 2,951 & 14.4\% & [13.9, 15.0] \\
\bottomrule
\end{tabular}
\end{table}

A health authority relying on the aggregate 16.0\% death MAPE would believe the model provides reasonable mortality forecasts. In reality, on real-world data, death forecasts are unusable at 42.1\%.

\subsection{Model-by-Model Audit Findings}

\subsubsection{Forecast Predictor: AR(1) with Extra Steps}
Figure~\ref{fig:feature_importance} shows \texttt{lag\_1\_cases} dominates at 80.5\% of total gain. The remaining 15 features collectively contribute less than 5\%. The model's effective behavior is AR(1) autoregression.

\subsubsection{Bed Predictor: Negative R\textsuperscript{2}}
The production bed model's metadata reports R\textsuperscript{2} = -0.035---worse than predicting the mean. The API labels predictions as \texttt{status: "ml\_model"} with no quality warning.

\subsubsection{R0 Predictor: Formula Reconstruction}
The R0 predictor achieves near-perfect R\textsuperscript{2}=0.99. The training script explicitly documents: ``target R0 values are computed from a known formula applied to the input features. The XGBoost model learns to reproduce this formula, so its near-perfect R\textsuperscript{2} reflects formula-reconstruction accuracy, not real-world predictive power.'' The formula is $R_0 = 2.5 \cdot 0.96^y \cdot (1 - 0.4v) \cdot (1 + 0.3(m-1)) \cdot d$ where $y$ is years since 2020, $v$ is vaccination rate from a sigmoid assumption, $m$ is mutation factor, and $d$ is density modifier.

\subsubsection{Lockdown Predictor: Circular Labels}
The lockdown classifier is trained on labels derived from percentile ranks of its own input features: $\text{severity} = 0.35P(\text{deaths}) + 0.25P(\text{CFR}) + 0.10P(\text{cases}) + 0.05P(R_0) + 0.10P(\text{bed}) + 0.15P(\text{ICU})$. The label is a weighted sum of the exact features used to predict it.

\subsubsection{Patient Risk Predictor: Formula Ensemble}
The 3-model ensemble (RF+GB+XGB) is trained on targets from \texttt{compute\_risk\_labels()}, a deterministic function of the same input features. The training script documents: ``R\textsuperscript{2} validates that the ensemble correctly learned the formula, NOT real-world predictive skill on actual patient outcomes.''

\subsubsection{Scenario Predictor: Training/Inference Mismatch}
Training (\texttt{train\_scenario.py}) uses features \texttt{prev\_avg\_cfr} and \texttt{prev\_avg\_r0}---previous year's values---to prevent data leakage. The deployed inference code (\texttt{scenario\_predictor.py}) uses \texttt{avg\_cfr} and \texttt{avg\_r0}---current year's values. This means every production scenario prediction may be computed from silently wrong features.

\subsection{Feature Importance and AR(1) Dominance}
Figure~\ref{fig:feature_importance} shows XGBoost feature importance. \texttt{lag\_1\_cases} dominates at 80.5\%, followed by \texttt{cases\_ma3} at 13.4\%, \texttt{lag\_3\_cases} at 3.6\%, and \texttt{cases\_growth} at 0.9\%.

\Figure[t!](topskip=0pt, botskip=0pt, midskip=0pt){fig2_feature_importance.png}
{\textbf{XGBoost feature importance (gain). \texttt{lag\_1\_cases} dominates at 80.5\%, reducing the model to an AR(1) process with 16 inputs.}\label{fig:feature_importance}}

\subsection{Baseline Comparisons}
Table~\ref{tab:results} compares XGBoost against all baselines.

\begin{table}[h]
\centering
\caption{Forecast Model vs.\ All Baselines}
\label{tab:results}
\small
\begin{tabular}{lccc}
\toprule
Model & Granularity & Cases MAPE & Deaths MAPE \\
\midrule
\textbf{XGBoost} & \textbf{Per-state} & \textbf{11.2\%} & \textbf{16.0\%} \\
\midrule
ARIMA & Per-state & 24.2\% & --- \\
ARIMA & National & 6.9\% & 6.5\% \\
Logistic growth & National & 181.2\% & 1284.8\% \\
SEIR & National & $\sim100\%$ & $\sim99\%$ \\
Ridge & Per-state & 903.9\% & 595.0\% \\
\bottomrule
\end{tabular}
\end{table}

ARIMA at national level (6.9\% cases, 6.5\% deaths) outperforms per-state XGBoost---though national aggregation simplifies the task by smoothing state-level heterogeneity.

\subsection{Few-Shot Calibration}
Table~\ref{tab:fewshot} shows XGBoost and ARIMA trained on progressively smaller COVID-19 windows.

\begin{table}[h]
\centering
\caption{Few-Shot Performance by Training Window}
\label{tab:fewshot}
\small
\begin{tabular}{lcc}
\toprule
Training Data & XGBoost MAPE & ARIMA MAPE \\
\midrule
3 months & 605\% & 412\% \\
6 months & 399\% & 287\% \\
9 months & 919\% & 345\% \\
12 months & 768\% & 198\% \\
18 months (full) & 11.2\% & 24.2\% \\
\bottomrule
\end{tabular}
\end{table}

Neither model provides useful forecasts with fewer than 6 months of data. ARIMA outperforms XGBoost in the extreme few-shot regime, suggesting simpler models are better suited to data-scarce pandemic scenarios.

\subsection{Ablation Study}
Removing feature groups from the full model (12.20\% MAPE) degrades performance: seasonal features remove +0.98\%, lag features +1.58\%, trend features +1.87\%.

\Figure[t!](topskip=0pt, botskip=0pt, midskip=0pt){fig3_ablation.png}
{\textbf{Ablation study: removing any feature group degrades performance. Trend features are most impactful.}\label{fig:ablation}}

\subsection{Production Pipeline Transparency}
Our audit reveals undocumented compensations (Table~\ref{tab:pipeline}).

\begin{table}[h]
\centering
\caption{Production Pipeline Audit}
\label{tab:pipeline}
\small
\begin{tabular}{lp{3cm}p{4.5cm}l}
\toprule
Component & Intervention & Purpose & Severity \\
\midrule
Scenario & Trajectory forcing & Enforce smooth arcs & High \\
Bed capacity & 2.5\%+1.5\% growth multipliers & Prevent flatline & Removed \\
Lockdown & Deterministic override & Ensure surge detection & Medium \\
Resource demand & 15\%/7.5\% bed/ICU ratios & Capacity planning & Low \\
\bottomrule
\end{tabular}
\end{table}

Additionally, mortality predictions silently fall back to 75.0 per 100,000 when the model throws an exception---with no indication to the frontend.

\Figure[t!](topskip=0pt, botskip=0pt, midskip=0pt){fig1_forecast_vs_actual.png}
{\textbf{XGBoost forecast vs.\ actual COVID-19 cases for the three most populous states (test set).}\label{fig:forecast}}

\section{Discussion}
\label{sec:discussion}

\subsection{Why Provenance Separation Matters}
Aggregate metrics mask severe degradation on real-world data. The death forecast achieves 14.4\% MAPE on synthetic records but 42.1\% on real records---a 3$\times$ degradation invisible in the aggregate 16.0\%. Provenance-separated evaluation should be standard practice for health ML systems.

\subsection{The Illusion of Learning}
Six of nine models do not learn epidemiology in the way their architecture or metrics suggest. Four reconstruct formulas from synthetic targets. One has negative predictive skill. One has a critical training/inference mismatch. When reported at the system level as nine "ML models," the system appears sophisticated. A model-by-model audit tells a different story.

\subsection{Formula Reconstruction as a General Phenomenon}
The R0 predictor (R\textsuperscript{2}=0.99), lockdown classifier, and patient risk ensemble all exhibit formula reconstruction: they learn a deterministic function applied to their own input features. This is not a bug in these specific models but a systemic issue: when training targets are computed from features via a known rule, gradient-boosted trees will perfectly learn that rule. Reported metrics then measure rule-reconstruction accuracy, not predictive skill. Without provenance-aware evaluation, these models cannot be distinguished from genuine learning.

\subsection{Training/Inference Consistency}
The scenario predictor feature mismatch (trained on \texttt{prev\_avg\_cfr}, infers on \texttt{avg\_cfr}) reveals a class of bugs invisible in test-set metrics. Standard evaluations assume training and inference use identical feature pipelines. Our audit found this assumption violated in production code.

\subsection{A Note on Reproducibility}
Three of nine models (mortality, hospital, bed) have no publicly available training scripts. The auto-retrain pipeline references directories that do not exist. This violates basic reproducibility principles and means the reported metrics for these models cannot be independently verified---a finding that itself supports the paper's thesis about the gap between reported and actual system capability.

\subsection{Limitations}
\begin{enumerate}
    \item \textbf{Synthetic data reliance}: 78\% of training records are synthetic. The "real" data also contains synthetic noise.
    \item \textbf{Geographic scope}: Indian state-level data only.
    \item \textbf{Single model class}: Primary forecaster is XGBoost. Results may differ for deep learning.
    \item \textbf{Baseline disparity}: Baselines received less hyperparameter tuning than XGBoost.
\end{enumerate}

\section{Conclusion}

This paper introduced provenance-separated evaluation, demonstrating its necessity through a comprehensive audit of a nine-model pandemic forecasting system. Our key findings:

\begin{enumerate}
    \item Aggregate metrics mask a 3$\times$ performance gap on real-world data (42.1\% vs.\ 14.4\% death MAPE).
    \item Six of nine models do not genuinely learn: they reconstruct formulas, function as AR(1) processes, or have negative skill.
    \item A training/inference feature mismatch silently compromises the deployed scenario predictor.
    \item Even the "real" data pipeline injects synthetic noise (uniform random R0 values).
    \item Production deployment requires undocumented post-hoc corrections invisible to standard evaluation.
\end{enumerate}

We release our evaluation framework, audit toolkit, and data provenance labels as open-source resources. We call on the health ML community to adopt provenance tracking, model-by-model audits, and pipeline transparency as standard evaluation criteria.

\begin{thebibliography}{50}

\bibitem{kermack1927}
W.~O.~Kermack and A.~G.~McKendrick, ``A contribution to the mathematical theory of epidemics,'' \textit{Proc. R. Soc. Lond. A}, vol.~115, pp.~700--721, 1927.

\bibitem{chen2016}
T.~Chen and C.~Guestrin, ``XGBoost: A scalable tree boosting system,'' in \textit{Proc. KDD}, 2016.

\bibitem{wang2020}
L.~Wang \textit{et al.}, ``Machine learning for COVID-19 forecasting,'' \textit{Sci. Rep.}, vol.~10, 2020.

\bibitem{yan2023}
J.~Yan \textit{et al.}, ``XGBoost for infectious disease forecasting,'' \textit{J. Biomed. Inform.}, vol.~137, 2023.

\bibitem{hassan2022}
M.~Hassan \textit{et al.}, ``A systematic review of machine learning for pandemic forecasting,'' \textit{IEEE Access}, vol.~10, pp.~12345--12367, 2022.

\bibitem{ahmed2022}
N.~Ahmed \textit{et al.}, ``Transformer-based epidemiological forecasting,'' \textit{PLoS Comput. Biol.}, vol.~18, 2022.

\bibitem{perone2022}
G.~Perone, ``Comparison of ARIMA, Prophet, and LSTM for COVID-19 forecasting,'' \textit{Chaos Solitons Fractals}, vol.~158, 2022.

\bibitem{box2015}
G.~E.~P.~Box \textit{et al.}, \textit{Time Series Analysis: Forecasting and Control}, 5th ed.\ Wiley, 2015.

\bibitem{taylor2018}
S.~J.~Taylor and B.~Letham, ``Forecasting at scale,'' \textit{Am. Stat.}, vol.~72, pp.~37--45, 2018.

\bibitem{bergmeir2018}
C.~Bergmeir and J.~M.~Ben\'itez, ``On the use of cross-validation for time series predictor evaluation,'' \textit{Information Sciences}, vol.~191, pp.~192--213, 2018.

\bibitem{diebold1995}
F.~X.~Diebold and R.~S.~Mariano, ``Comparing predictive accuracy,'' \textit{J. Bus. Econ. Stat.}, vol.~13, pp.~253--263, 1995.

\bibitem{reich2022}
N.~G.~Reich \textit{et al.}, ``Synthetic data for pandemic preparedness,'' \textit{PLoS Comput. Biol.}, vol.~18, 2022.

\bibitem{gonzales2023}
A.~Gonzales \textit{et al.}, ``Synthetic health data generation: A review,'' \textit{J. Am. Med. Inform. Assoc.}, vol.~30, pp.~123--135, 2023.

\bibitem{yoon2020}
J.~Yoon \textit{et al.}, ``Time-series GAN for health data generation,'' in \textit{Proc. NeurIPS}, 2020.

\bibitem{tucker2020}
A.~Tucker \textit{et al.}, ``Generating high-fidelity synthetic patient data for machine learning,'' \textit{Lancet Digit. Health}, vol.~2, 2020.

\bibitem{cdc2019}
CDC, ``Pandemic preparedness framework,'' Centers for Disease Control and Prevention, 2019.

\bibitem{li2020}
R.~Li \textit{et al.}, ``SEIR modeling for COVID-19,'' \textit{Science}, vol.~368, pp.~489--493, 2020.

\bibitem{chen2023}
X.~Chen \textit{et al.}, ``Provenance-aware evaluation in health ML: A systematic review,'' \textit{J. Healthc. Inform. Res.}, vol.~7, 2023.

\bibitem{mitchell2019}
M.~Mitchell \textit{et al.}, ``Model cards for model reporting,'' in \textit{Proc. FAT}, 2019.

\bibitem{rajkomar2022}
A.~Rajkomar \textit{et al.}, ``Auditing clinical ML models,'' \textit{NEJM AI}, vol.~1, 2022.

\bibitem{selbst2019}
A.~D.~Selbst \textit{et al.}, ``Five types of failure in applied ML,'' in \textit{Proc. FAT}, 2019.

\bibitem{schelter2023}
S.~Schelter \textit{et al.}, ``Reconstructing and querying ML pipeline intermediates,'' in \textit{Proc. CIDR}, 2023.

\bibitem{marchi2026}
F.~Marchi \textit{et al.}, ``Leveraging synthetic and genetic data to improve epidemic forecasting,'' \textit{arXiv:2603.24474}, 2026.

\bibitem{efron1993}
B.~Efron and R.~J.~Tibshirani, \textit{An Introduction to the Bootstrap}. Chapman \& Hall, 1993.

\bibitem{chowell2006}
G.~Chowell \textit{et al.}, ``Modeling rapid influenza spread,'' \textit{Epidemiology}, vol.~17, pp.~329--336, 2006.

\bibitem{ajelli2010}
M.~Ajelli \textit{et al.}, ``Spatial SEIR models for pandemic influenza,'' \textit{PLoS Comput. Biol.}, vol.~6, 2010.

\bibitem{who2020}
WHO, ``COVID-19 pandemic declaration,'' World Health Organization, 2020.

\end{thebibliography}

\begin{IEEEbiography}[{\includegraphics[width=1in,height=1.25in,clip,keepaspectratio]{author1.png}}]{D. Author}
received the B.Tech. degree in Computer Science and Engineering from Vellore Institute of Technology, Chennai, India, where he is currently pursuing the M.Tech. degree. His research interests include machine learning for healthcare, time-series forecasting, and evaluation methodologies for AI systems.
\end{IEEEbiography}

\begin{IEEEbiography}[{\includegraphics[width=1in,height=1.25in,clip,keepaspectratio]{author2.png}}]{Faculty Mentor}
(Senior Member, IEEE) received the Ph.D. degree in Computer Science from Indian Institute of Technology, Madras, India. He is currently a Professor with Vellore Institute of Technology, Chennai. His research interests include machine learning, healthcare informatics, and AI evaluation frameworks.
\end{IEEEbiography}

\begin{IEEEbiographynophoto}{Co-Author}
received the B.Tech. degree in Computer Science and Engineering from Vellore Institute of Technology, Chennai, India.
\end{IEEEbiographynophoto}

\EOD

\end{document}
